% ML4H 2026 Findings track submission (BIOP02)
% !! 제출 전 교체 필요 !!
% ML4H 2026 공식 템플릿은 Overleaf 전용이라 이 파일에 포함할 수 없다.
%   https://www.overleaf.com/latex/templates/machine-learning-for-health-ml4h-2026-template/sqgwhtyswgcy
% 아래 \documentclass ~ \begin{document} 구간을 템플릿 프리앰블로 교체하고
% 본문은 그대로 옮기면 된다. 본문은 템플릿 의존 명령을 쓰지 않는다.
% 이중맹검: 저자·소속·사사는 camera-ready 에만 넣는다.
\documentclass[10pt]{article}
\usepackage[a4paper,margin=2.2cm]{geometry}
\usepackage{graphicx}\usepackage{amsmath,amssymb}
\usepackage[T1]{fontenc}\usepackage{lmodern}
\usepackage[hidelinks]{hyperref}\usepackage{caption}
\setlength{\parskip}{4pt}\setlength{\parindent}{0pt}
\title{Safety of Substitution: Cost-Aware Evaluation of H\&E-Based Molecular Test Replacement Across Five Cancers}
\author{}\date{}
\begin{document}
\maketitle
\texttt{}`markdown

\begin{abstract}
That a tumour's molecular phenotype can be predicted from haematoxylin-and-eosin (H\&E) histology, and that such a prediction may clinically replace molecular testing, are different claims. We propose a cost-of-substitution frame that converts prediction errors into the misassignment cost of a pre-specified treatment routing. Across five cancers (breast as the anchor plus lung, colorectal, gastric and head and neck), we sealed and tested a pre-registered morphological-correlate law under one protocol, but powered confirmation is limited. The observed spectrum is fixed by measurement at two ends: head and neck HPV, a morphology-legible viral axis, reached holdout AUROC 0.959 and was the only powered, non-control positive result in the pre-registered holdout of the primary model UNI; lung LUSC histology reached 0.939 as a positive control. By contrast, breast HER2 showed no signal supporting substitution in this cohort and routing definition, with AUROC 0.599 and anti-HER2 misassignment rate 1.00. The underlying HER2 phenotype prediction remained near chance under H\&E stain normalisation, so the HER2 negative is not an artefact of stain variation, although the routing/cost step itself was not re-run under normalisation. Most clinically important mutation and amplification axes are exploratory because holdout positives fall short of 25 under our pre-registered site-disjoint split. Some exploratory results are compatible with the morphological-correlate hypothesis, but we lacked the power to establish the direction of the law across mutation and amplification axes. We propose safety of substitution, not predictability, as the decision criterion for this retrospective, cohort-level map.
\end{abstract}
\section*{1. Introduction}

AI analysis of histopathological H\&E images has expanded across organs as digital pathology has spread [CITE-I1], with weakly supervised multiple-instance learning and pathology foundation models supporting work in urological cancers, breast cancer, pancreatic cancer and other settings [CITE-I2] [CITE-I3] [CITE-I4] [CITE-I5] [CITE-I6]. A persistent target has been molecular-state prediction from images, because IHC staining and tissue-destructive molecular tests are generally costly and slow, whereas H\&E is relatively inexpensive and already acquired in routine care [CITE-I7]. Molecular tests guide diagnosis, prognosis and treatment across cancer types [CITE-I8], and molecular state can often be predicted from H\&E [CITE-I9].

Prediction alone, however, does not establish that a molecular test can be replaced. The same AUROC can have different clinical consequences depending on which treatment decision the error changes [CITE-I10] [CITE-I11]. We therefore ask when H\&E substitution is safe, not simply when a label is predictable. Our frame converts prediction errors into the misassignment cost of treatment routing. It does not predict drug response and takes no drug structure as input. It operationalises only the substitution cost from marker to treatment assignment.

We test a pre-registered morphological-correlate law across five cancers: breast as the anchor plus lung, colorectal, gastric and head and neck. The law states that H\&E can cheaply stand in for a test only when the molecular alteration has a morphological correlate recognisable at H\&E resolution. These cancers are a deliberate boundary, not an open pan-cancer atlas expansion; and sealing predictions before results provides claim discipline rather than automatic confirmatory strength.

This paper introduces a cost-of-substitution frame, applies one pre-registered protocol across five cancers, reports an honest negative anchor for breast HER2, adjudicates insufficient power on mutation and amplification axes rather than reporting only high-scoring axes, and frames a different question from single-cohort breast prediction [CITE-I12], including Fernandez-Romero 2026 [CITE-I12], and from drug-sensitivity prediction [CITE-I13]: when substitution for a molecular test is safe.

\section*{2. Results}

We tested about fifteen endpoints across five cancers under pre-registration. Exactly one powered, non-control confirmation was obtained: HPV in head and neck cancer. Most remaining mutation and amplification axes fell short of twenty-five holdout positives and therefore remain undecided under the pre-registered rule. We therefore do not claim that the law was validated across five cancers. We establish two things: both ends of the map are fixed by measurement under one protocol, and the middle of the map cannot yet be decided because of the power ceiling described in R2.

\subsection*{R1. Substitution-Cost Spectrum}

This section converts per-axis substitutability into cost. Figure 1 overlays cost on the confusion matrix weighted by therapeutic distance; per-axis cost and the confidence interval of the headline contrast are shown in Figure A1.

In the holdout of the primary model UNI [CITE-M5], head and neck HPV reached AUROC 0.959 [0.921--0.986] with 26 positives, above the pre-registered threshold of 0.80. This axis is shaped by viral infection rather than by a mutation, and may extend the law's morphological-correlate clause to non-keratinising, basaloid morphology, but confirmation is confined to HPV. Positive controls behaved as expected: lung LUSC histology 0.939 [0.905--0.967] with 153 positives, and head and neck grade 0.815 [0.742--0.882] with 41 positives.

At the illegible end, breast HER2 was 0.599, effectively at chance; in this cohort and routing definition it shows no signal supporting H\&E-based substitution and serves as the negative anchor. Gastric ERBB2 amplification was 0.644, while the shuffle-null was 0.641, so no signal should be inferred. Lung KRAS-G12C was 0.681, but a baseline using no image and predicting from histology alone reached 0.793; the apparent predictive power therefore comes from LUAD skew rather than mutation morphology.

Every endpoint is reported alongside a shuffle-null, a prevalence baseline and a pixel-mean baseline; lung mutation axes also carry a subtype-only baseline. Epistemic status differs by cohort: lung, gastric and head and neck were sealed-forward tests with predictions committed before results, whereas colorectal was analysed after results were already available and is excluded from the tally of powered, sealed confirmations.

\textbf{Table R1. Observed substitution-cost spectrum (UNI canonical). Cost is treatment misassignment loss.}

\begin{table}[ht]\centering\footnotesize
\begin{tabular}{lllllll}\hline
Cancer & Axis & Role & AUROC [95\% CI] & Holdout n\_pos / control baseline & Morphological correlate & Verdict \\\hline
Head and neck & HPV & Legible axis (viral) & 0.959 [0.921--0.986] & 26 positives & Present (non-keratinising, basaloid) & Single-FM, site-disjoint confirmation (model-independence and site confounding untested) \\
Lung & LUSC histology & Positive control & 0.939 [0.905--0.967] & 153 positives & Morphology itself & Passes chance-exclusion; excluded from the confirmation tally (site signature inseparable) \\
Head and neck & Grade & Positive control & 0.815 [0.742--0.882] & 41 positives & Present & Pass \\
Colorectal & BRAF V600E & Retrospective & 0.882 [0.817--0.938] & 15 positives & Present (serrated/MSI co-occurrence) & Consistent; retrospective, underpowered, exploratory \\
Gastric & MSI-H & Legible axis & 0.860 (development 0.899) & 24 positives & Present (immune) & Undecided (1 short) \\
Lung & EGFR activating & Graded & 0.852 & 15 positives & Partial & Undecided \\
Lung & KRAS-G12C & Required axis & 0.681 (subtype-only 0.793) & 14 positives & Absent (histology skew) & Undecided \\
Gastric & ERBB2 amplification & Required axis & 0.644 (shuffle 0.641) & 14 positives & Absent & Underpowered; no observed signal \\
Breast & HER2 & Anchor & 0.599 & near-random & Absent & No signal supporting substitution; negative anchor \\
Gastric & Lauren diffuse & Original positive control & 0.536 (development 0.963) & pixel-mean 0.631 & Weakly present & Site-split case (R4) \\
\hline\end{tabular}\end{table}

HPV robustness passed 5-seed chance-exclusion in only 2 of 3 foundation models: UNI and UNI2-h [CITE-M23] passed, while Virchow2 [CITE-M6] did not clear the pre-specified criterion (real 0.9199 < threshold 0.9234, margin $-$0.0035). The site audit also found site-label structuring (Cram\'er's V = 0.397). HPV is therefore the single powered anchor that fixes one end of the map, not a generalisation of the law and not a model-independent confirmation. For lung histology, V(site, label) = 1.000, so morphology and site signature cannot be separated. The HPV and MSI-H axes were robust to tile subsampling, but slide-level reproducibility was not established; predictions diverged between multiple slides from the same patient (TCGA-QK-A6IF, 0.92 vs 0.001). Clinical substitution on these axes therefore requires multi-centre prospective validation, multi-slide reproducibility and verification of tumour-region contribution.

\subsection*{R2. Power Ceiling}

Clinically important mutation and amplification axes repeatedly fell short of the threshold of twenty-five positives in the pre-registered holdout. In the single site-disjoint split we chose, most actionable mutations were not adequately powered. We do not generalise this to public data being impossible in principle; power might be recovered with grouped or leave-one-site-out cross-validation, which we treat as exploratory.

The affected axes were lung EGFR activating, 15 positives; lung KRAS-G12C, 14 positives; gastric ERBB2 amplification, 14 positives; gastric MSI-H, 24 positives; gastric EBV, 7 positives; and head and neck EGFR amplification, 17 positives; full verdicts are in Appendix A, Table R2. The threshold was not adjusted after the fact. Gastric MSI came one patient short at 24, and we did not lower the criterion from 25 to 24. Deciding substitutability of mutation axes will therefore require institutional cohorts or prospective collection, and until then the middle of the map is left open.

\subsection*{R3. Breast Anchor}

Under the pre-specified routing definition and in this cohort, anti-HER2 assignment based on H\&E-predicted subtype did not support identification of treatment candidates: anti-HER2 misassignment rate 1.00. The misassignment rate and cost interpretation depend on the stated operating point, and the analysis that pre-specifies that threshold is not yet complete. This is an observation showing, as misassignment loss, that H\&E substitution may not be safe in this region.

Per-axis cost flips between endocrine therapy and chemotherapy depending on the routing scheme (0.378 versus 0.035; 0.105 versus 0.510). The only claims robust to a change of scheme are the anti-HER2 misassignment rate of 1.00 and the fact that the confidence interval of the headline contrast excludes 0; we do not extend these into a claim that other axes are safe.

The HER2 negative is not an artefact of stain variation at the level of phenotype prediction. In a stain-normalisation robustness check on the breast anchor, using Macenko normalisation with embeddings re-extracted and CLAM attention MIL [CITE-M9] re-trained on the same folds, HER2 phenotype prediction remained near chance: AUROC 0.641 versus 0.599 without normalisation in Table R1. ER stayed high (0.917 versus 0.901) and PAM50 was preserved (0.740 versus 0.759). The rank of the anchor endpoints, ER high > PAM50 mid > HER2 near-chance, was preserved with and without stain normalisation. Two limits remain: the routing/cost pipeline was not re-run under normalisation, and no shuffle-null was computed for stain-normalised runs. This robustness check covers the breast anchor only.

\subsection*{R4. Gastric Lauren}

Lauren diffuse was originally a positive control. Signet-ring and diffuse-type tumours have strong H\&E morphology and should have scored high, yet the result was 0.536. The cause is not illegibility. In the same pipeline, gastric MSI generalised from 0.899 in development to 0.860 in holdout, whereas Lauren fell from 0.963 to 0.536, a drop of 0.43. The low-resolution pixel-mean baseline, 0.631, exceeded the MIL model, 0.536, so a weak morphological signal exists and the model failed to capture it. The direct cause is that Lauren prevalence varies sharply across institutions and the site-disjoint split concentrated high-prevalence institutions in the evaluation set: 46 per cent in training versus 88 per cent in evaluation.

We therefore describe this case as a methodological instance in which site-disjoint evaluation correctly blocked shortcut learning, and confine the low-confidence verdict to gastric Lauren. MSI in the same cohort remains valid. We do not write that H\&E cannot see Lauren. The representative cases for the thesis that predictability and substitutability are different claims are breast HER2 and lung KRAS.

\subsection*{R5-R7. Appendix-Supported Checks}

R5 foundation-model robustness is in Appendix A, Table R5. In brief, lung endpoint ordering was preserved across UNI, Virchow2 and UNI2-h (histology > EGFR > KRAS; Spearman 1.000 against UNI for both newer models), and the principal negatives reproduced. However, individual molecular axes did not clear chance-exclusion in all models: HPV passed in UNI and UNI2-h but not Virchow2, and colorectal BRAF passed in UNI and Virchow2 but not UNI2-h. This is rank stability and negative-result reproducibility, not model independence of the law.

R6 external treatment-outcome anchoring is retained only as a pending pointer, with method and provisional value in Appendix C. The anti-HER2 axis score was computed by frozen transfer and used to stratify pathological complete response in an external cohort, evaluated by AUROC with bootstrap confidence intervals and compared against a measured-HER2 probability baseline with DeLong's test [CITE-M12]. This result is \texttt{critic\_status: pending} and is not promoted to the Abstract or headline claims. Directionally, the anti-HER2 axis did not stratify pCR from the H\&E-predicted phenotype, consistent with the retrospective map's HER2 negative.

R7 spatial transcriptomics analyses are in Appendix B because they are \texttt{hypothesis\_only} and Critic pending. In brief, HER2-positive tumours showed low-expression tumour regions that may help explain subtype-routing error, but mRNA differs from protein and amplification, a spot is not a cell, and the spatial cohort is not the same cohort. The colorectal spatial correlate did not emerge at Visium resolution, leaving that mechanism open.

\section*{3. Discussion}

Our map takes safety of substitution, not predictability, as its criterion. The frame flags axes where H\&E substitution is dangerous, such as breast HER2 and lung KRAS, and separates them from axes that are undecided or artefactual. Gastric Lauren is excluded from the dangerous-substitution list because it is a site-split artefact rather than an absence of morphology.

In breast HER2, routing from predicted subtype failed consistently in this cohort and routing definition, indicating as cost a region where molecular testing remains necessary. We state the scheme dependence of per-axis cost and restrict robust claims to the anti-HER2 misassignment rate of 1.00 and the confidence interval of the contrast. This negative is not a stain artefact at the phenotype-prediction level; the external treatment-outcome anchor is directionally consistent with it, but remains \texttt{critic\_status: pending}.

Fernandez-Romero 2026 [CITE-I12] is convergent prior work on the HER2 axis rather than a priority boundary. The difference is the question asked: they quantify how much performance drops across breast cohorts, whereas we ask when substitution for a molecular test is safe. The output metric also differs: they report macro-F1/PR-AUC degradation, while we report misassignment rate against a pre-defined treatment-routing rule. The scope differs as well: their study is breast-only, whereas ours uses breast as an anchor plus lung, colorectal, gastric and head and neck. Finally, their split design is patient-stratified random cross-validation, while ours is pre-registered site-disjoint evaluation with sealed scoring. Any causal explanation for the difference in degradation is only a hypothesis, because metrics, foundation models and inclusion criteria differ.

The limits are placed in front. All results are retrospective, cohort-level and \texttt{hypothesis\_only}; they are not claims of individual-level benefit. The site-disjoint split prevented leakage in which slides from the same tissue source site enter both training and evaluation, but label-institution coupling remains and limits interpretation as pure morphology. The audit confirmed site-label structuring in multiple endpoints already described above, including lung histology at V = 1.000 and head and neck HPV at V = 0.397. This is a necessary condition for confounding, not proof that the model reads site. Site predictability from H\&E and leave-one-site-out performance are needed to adjudicate the confounding question.

Stain variation is a separate limit. H\&E stain normalisation was not applied in the main pipeline; uncorrected stain variation is a known source of domain shift in pathology imaging [CITE-M17]. The breast-anchor robustness check preserved the HER2-negative pattern, but the cross-cancer headline axes, head and neck HPV and lung histology, were not stain-verified. The two headline results most vulnerable to scanner or stain critique therefore carry site-confounding flags and are not yet stain-verified.

On model independence, the relative AUROC ordering among lung endpoints and the principal negative results were preserved across three foundation models, whereas whether an individual molecular axis clears chance-exclusion varied by model. This is ordering stability, not model independence of the law as a whole.

PAM50 provenance is stated explicitly because the endpoint carries weight in the breast anchor. The canonical PAM50 endpoint is CLAM-MB, uni\_v1, 4-class (\texttt{pam50\_clam\_mb\_uni\_v1\_4class}), not the v2 CLAM-SB endpoint used for ER/PR/HER2. PAM50 labels from the manifest nearest-centroid computation [CITE-M2] show 57.0 \% concordance (514/902) against cBioPortal PanCancer Atlas SUBTYPE labels [CITE-M3], i.e. 43.0 \% discordance (388/902), with largest disagreements LumB$\leftrightarrow$LumA and Normal$\rightarrow$LumA. Coverage was 97.2 \%, so the fallback condition for using local labels because cBioPortal coverage is short was not met. This remains a label-source reconciliation limit, not a hidden error.

Clinically, this observational map identifies negative axes where H\&E substitution is clearly dangerous, undecided axes that present data cannot adjudicate, and morphology-legible axes that justify prospective validation. The paper makes no clinical recommendation and no claim of wholesale replacement. Its contribution is not beating the gold standard, but making predictable, as a map, when inexpensive H\&E can pre-screen or triage molecular testing and when it cannot.

\section*{4. M5. Cost-of-Substitution Frame}

Substitution cost is defined by multiplying the confusion matrix by therapeutic distance, giving the misassignment cost incurred where the treatment chosen from the measured marker and the treatment chosen from the H\&E-predicted marker diverge. The lead indicator is the distance-independent misroute rate. This frame does not predict drug response and takes no drug structure as input. Cohorts and labels (M1), evaluation design (M4) and claim discipline (M6) are moved in full to Appendix C.

\section*{Appendix A. Results Tables and R5 Foundation-Model Robustness}

\textbf{Table R2. Power ceiling}

\begin{table}[ht]\centering\footnotesize
\begin{tabular}{lll}\hline
Axis & Holdout positives & Verdict \\\hline
Lung EGFR activating & 15 & Undecided \\
Lung KRAS-G12C & 14 & Undecided \\
Gastric ERBB2 amplification & 14 & Underpowered; no observed signal \\
Gastric MSI-H & 24 & Undecided (1 short of threshold) \\
Gastric EBV & 7 & Exploratory \\
Head and neck EGFR amplification & 17 & Undecided \\
\hline\end{tabular}\end{table}

Holding slides, site-disjoint holdout and endpoints fixed, we retrained CLAM after swapping only the embedding space: UNI 1024-d, Virchow2 2560-d and UNI2-h 1536-d. The claim concerns ordering, not absolute values, and not model independence of individual axes.

Across all three embedding spaces, lung endpoints kept the order histology > EGFR > KRAS. Spearman correlation against UNI was 1.000 for both newer models, and 5-seed chance-exclusion passed 6 of 6 in lung. This is ordering stability, not model generality and not confirmation.

Single-endpoint results diverged by model. Head and neck HPV passed in UNI and UNI2-h, but in Virchow2 the point estimate was high (0.9199) while the shuffle-null spread was wider, so it did not clear the pre-specified criterion. Colorectal BRAF passed in only two of three models: UNI and Virchow2; UNI2-h did not clear the low-power shuffle-null.

Two kinds of negative must be distinguished. Gastric Lauren fails chance-exclusion in all three (0.536 / 0.640 / 0.603), reproducing the site-confounding failure of R4 and not proving absence of morphological signal. Gastric ERBB2 amplification fails in all three (0.644 / 0.668 / 0.585) because real $\approx$ null in every case, consistent with absence of signal. Lumping them would re-introduce the error R4 corrects.

Head and neck EGFR amplification shows that clearing chance-exclusion is not itself evidence of signal. It is formally recorded as passing in UNI2-h, yet its real AUROC is 0.505, essentially chance; the shuffle-null spread was narrow, so the threshold sat correspondingly low.

\textbf{Table R5. Multiple foundation models (5-seed canonical)}

\begin{table}[ht]\centering\footnotesize
\begin{tabular}{lllll}\hline
Endpoint & UNI & Virchow2 & UNI2-h & 5-seed chance-exclusion \\\hline
Head and neck HPV & 0.9594 & 0.9199 & 0.9559 & UNI and UNI2-h pass / Virchow2 fails (margin $-$0.0035) \\
Colorectal BRAF & 0.8676 & 0.8798 & 0.8978 & UNI and Virchow2 pass / UNI2-h fails \\
Gastric MSI-H & 0.8599 & 0.8795 & 0.8670 & Both newer models pass \\
Gastric Lauren (positive control) & 0.5364 & 0.6404 & 0.6033 & All three fail (site-confounding, reproduced) \\
Gastric ERBB2 amplification & 0.6444 & 0.6682 & 0.5845 & All three fail (signal absent, reproduced) \\
Lung (histology > EGFR > KRAS) & Order preserved & Order preserved & Order preserved & Spearman 1.000; 6/6 \\
\hline\end{tabular}\end{table}

\section*{Appendix B. R7 Spatial Transcriptomics}

This exploratory analysis examines a possible mechanism behind the map using public spatial transcriptomics. Everything here is \texttt{hypothesis\_only} and has not passed Critic, so it should be read as mechanistic support rather than as a headline.

Even in confirmed HER2-positive tumours (8 patients), some tumour spots show ERBB2 levels indistinguishable from the non-tumour reference on the same section. The median probability that a tumour spot's ERBB2 falls below the reference is 0.158; in all 8 patients the confidence interval excludes 0 and the kill-test that rules out diffusion and depth artefacts is passed (interior-only 7/8, depth-conditioned 3/3). Because a patient carries only one label, these low-expression regions cannot be represented. This strengthens a candidate mechanism for subtype-routing error: HER2 may not be substitutable because the label discards information, rather than because the prediction is noisy. We do not assert a limit in principle. Limits are that mRNA differs from protein and amplification, a spot is not a cell, and the ST cohort is not our TCGA cohort.

By contrast, the spatial correlate predicted in colorectal did not emerge at Visium resolution. A 55 $\mu$m spot is coarser than nuclear resolution and cannot reach lymphocyte-specific texture; this is a substrate and resolution limit, not a biological refutation. The colorectal spatial mechanism remains open, and the appropriate test is a substrate with co-registered H\&E.

\section*{Appendix C. Additional Methods}

\subsection*{M1. Cohorts and Labels}

Breast cancer (TCGA-BRCA, about 1,010 diagnostic slides) [CITE-M1] served as the anchor, together with lung (TCGA-LUAD/LUSC) [CITE-M18], colorectal (TCGA-COAD/READ) [CITE-M19], gastric (TCGA-STAD) [CITE-M20] and head and neck (TCGA-HNSC) [CITE-M21], five cancers in total. Slide counts per cohort were measured in result JSON files: colorectal 523, lung 1,026, gastric 439, head and neck 468.

For the breast PAM50 endpoint, the canonical analysis endpoint is CLAM-MB, uni\_v1, 4-class (\texttt{pam50\_clam\_mb\_uni\_v1\_4class}). The labels use a nearest-centroid computation [CITE-M2] as implemented in genefu [CITE-M22]. The manifest labels and cBioPortal PanCancer Atlas SUBTYPE labels [CITE-M3] agree on 57.0 \% of overlapping patients (514/902; 43.0 \% discordance). The discordance is reported transparently.

\subsection*{M2. Tiling and Embedding}

Each whole-slide image was tiled into 256$\times$256 pixel patches at 20$\times$ magnification, tissue was separated from background by Otsu thresholding [CITE-M4], and a cap of 5,000 tiles per patient was imposed. The headline embedding is UNI v1 (1024-d) [CITE-M5]. For the model-independence test, the same coordinates were re-extracted with Virchow2 [CITE-M6] (2560-d, CLS token concatenated with mean patch token, register tokens excluded) and UNI2-h (1536-d) [CITE-M23]. The slide-level EXAONE Path 2.0 [CITE-M7] interface is incompatible with the coordinate-based pipeline and was excluded from the robustness set. Tiles were resized to 224$\times$224 and channel-normalised with ImageNet [CITE-M8] training-set statistics, using the standard torchvision constants. H\&E stain normalisation was not applied in the main pipeline.

\subsection*{M3. Model and Training}

We used CLAM-SB attention MIL [CITE-M9] with hidden 512, attention 256, 40--50 epochs and seed fixed at 42. Predictions were produced per slide and then aggregated per patient.

\subsection*{M4. Evaluation Design}

All evaluation was performed on a site-disjoint holdout. Slides from the same tissue source site were prevented from entering training and evaluation simultaneously, blocking leakage via institutional fingerprints [CITE-M10]; validation and test were combined for power. Three controls were used: a shuffle-null, a prevalence baseline (0.5) and a subtype-only or pixel-mean baseline. Confidence intervals are reported as 1,000-fold bootstrap 95\% CIs [CITE-M11]; where patient clustering matters, CIs were recomputed at the patient level.

\subsection*{M6. Claim Discipline}

Adjudication thresholds are cited only from the sealed pre-registration document, not from slides or observed values. The power rule, fewer than 25 positives $\rightarrow$ exploratory $\rightarrow$ INCONCLUSIVE, is not moved after seeing results and is applied symmetrically to confirmation and refutation. All outputs are \texttt{hypothesis\_only} and retrospective.

\subsection*{M7. External pCR Anchor}

The anti-HER2 axis score was computed by frozen transfer, meaning the anchor model was applied without further training, used to stratify pCR in an external cohort [CITE-M25], and evaluated by AUROC with bootstrap 95\% confidence intervals. It was then compared with the measured-HER2 probability baseline using DeLong's test [CITE-M12]. The pre-specified benchmark was defined as approaching and overlapping the benchmark of Farahmand and colleagues [CITE-M13] (0.80, 95\% CI 0.69 to 0.88). This anchor is provisional and is not carried into the Abstract or headline claims.

\subsection*{M8. Multi-Model Robustness}

Because embedding spaces are not interchangeable across foundation models [CITE-M14], CLAM was refitted from scratch in each space so that the comparison is made at the same level. The adjudication criterion is 5-seed shuffle-null chance-exclusion: real AUROC > null mean + 2 $\times$ standard deviation, ddof = 1, with seeds 42, 1, 2, 3 and 4. Determinism was verified by 2 re-runs at identical seeds. Final multi-FM Critic sign-off is in progress.

\subsection*{M9. Site/Batch Confounding Audit}

For each endpoint, we quantified whether the site-disjoint split confounds the label with the tissue source site. We computed Cram\'er's V [CITE-M15] between site and label with a permutation p-value, the train/test prevalence shift, and a permutation test of the site concentration of test positives. This analysis examines a necessary condition for confounding; final adjudication of whether the model actually uses site rests on site predictability from H\&E and on leave-one-site-out performance.

\subsection*{M10. Stain-Normalisation Robustness}

To test whether the anchor results are an artefact of uncorrected H\&E stain variation, embeddings were re-extracted from the breast-anchor slides with Macenko stain normalisation [CITE-M16], torchstain 1.3.0 [CITE-M24], and a fixed dense-tissue reference tile. CLAM was re-trained on the same folds for ER, HER2 and PAM50. Only phenotype prediction was re-run, not the routing/cost pipeline, and no shuffle-null was computed for the stain-normalised runs. This robustness check covers the breast anchor only; cross-cancer re-extraction under stain normalisation is deferred.

\end{document}
